% SEGMENT OLP-0591-S01
\documentclass[../../../include/open-logic-section]{subfiles}

\begin{document}

\olfileid{sth}{card-arithmetic}{fix}
\olsection{$\aleph$-நிலைப்புள்ளிகள்}

\olref[spine][]{chap} இல், மாற்றீட்டு அடிகோளத்துக்கு ஒரு
நியாயப்படுத்தல் தேவை எனக் குறிப்பிட்டோம். ஏனெனில் அது கணப்
படிநிலைமுறையை மிகவும் உயரமாக ஆக்குகிறது. கண அளவெண் கணிதத்தைச்
சிறிது கற்றுள்ள இப்போது, அந்த உயரத்தை ஓரளவு விளக்கலாம்.

$0 < \aleph_0$, $1 < \aleph_1$, $2 < \aleph_2$\ldots
என்பவை தெளிவு; ஒவ்வொரு படியிலும் அளவுவேறுபாடு உண்மையில்
\emph{பெரிதாகவே} செல்கிறது. எனவே எல்லா வரிசையெண் $\kappa$
க்கும் $\kappa< \aleph_\kappa$ என்று ஊகிக்கத் தோன்றலாம்.

ஆனால் $\ZFC$ ஐ ஏற்றால் இந்த ஊகம் \emph{தவறு}. உண்மையில்,
$\kappa=\aleph_\kappa$ என நிறைவேற்றும் கண அளவெண்களான $\kappa$,
\emph{$\aleph$-நிலைப்புள்ளிகள்} இருப்பதை நிறுவலாம்.

\begin{prop}\ollabel{alephfixed}
ஓர் $\aleph$-நிலைப்புள்ளி உள்ளது.
\end{prop}

\begin{proof}
மீள்வரையறையைப் பயன்படுத்தி பின்வருவனவற்றை வரையறுக்கவும்:
\begin{align*}
	\kappa_0 &= 0\\
	\kappa_{n+1} &= \aleph_{\kappa_n}\\
	\kappa&= \bigcup_{n < \omega}\kappa_n
\end{align*}
\olref[cardinals][classing]{unioncardinalscardinal} இன்படி
$\kappa$ ஒரு கண அளவெண். மேலும்,
\[
	\kappa= \bigcup_{n < \omega} \kappa_{n+1} =
	\bigcup_{n < \omega}\aleph_{\kappa_n} =
	\bigcup_{\alpha < \kappa}\aleph_\alpha = \aleph_\kappa
\]
% கட்டமைப்பால், $\kappa$ என்பது ஒவ்வொரு $\kappa_n$ ஐவிடவும் பெரிய
% மீச்சிறிய கண அளவெண். ஆகவே $\aleph_\kappa$ என்பது ஒவ்வொரு
% $\aleph_{\kappa_n}$ ஐவிடவும், அதனால் ஒவ்வொரு
% $\kappa_{n+1}$ ஐவிடவும் பெரிய மீச்சிறிய கண அளவெண். அதேபோல்
% $\kappa$ என்பதும் ஒவ்வொரு $\kappa_{n+1} = \aleph_{\kappa_n}$
% ஐவிடவும் பெரிய மீச்சிறிய கண அளவெண். எனவே $\kappa=
% \aleph_\kappa$.
\end{proof}

% SEGMENT OLP-0591-S02
நாம் இப்போது கட்டிய $\aleph$-நிலைப்புள்ளியைப் பற்றியே
பூலோஸ் ஒரு கட்டுரை எழுதினார். அதன் தொடக்கத்தில் $\kappa$
இருப்பதைக் குறிப்பிட்ட பிறகு அவர் சொன்னது:
\begin{quote}
	[கணக் கோட்பாட்டை முன்பு அறியாதவர்களின் பார்வையில் $\kappa$]
	\emph{மிகப் பெரிய} எண்; குறைந்தது $\kappa$ பொருள்கள்
	இருக்கின்றன என்று கூறும் எந்தக் கோட்பாட்டின் உண்மையையும்
	கேள்விக்குள்ளாக்கும் அளவுக்கு அது பெரியதாக எனக்குத் தோன்றுகிறது.
	\citep[p.~257]{Boolos2000}
\end{quote}
கட்டுரையின் முடிவில் அவர் கேட்ட கேள்வி:
\begin{quote}
	[கதையின்] தொடக்கத்தில் நிலைமை எப்படியிருந்தாலும்,
	இவ்வளவு தொலைவு வந்தபின் சக்கரங்கள் வெறுமனே சுழல்கின்றனவா?
	உண்மையில் உள்ள எதையும் விவரிப்பதை நாம் இனி கேட்கவில்லை
	என்று ஐயுற வேண்டுமா?
	\citep[p.~268]{Boolos2000}
\end{quote}
உண்மையில் உள்ள ``எதையும்'' நாம் கடந்து சென்றிருந்தால்,
அதற்கான பொறுப்பை நேராக மாற்றீட்டு அடிகோளத்திடமே சுமத்த
வேண்டும். ஏனெனில் அந்த அடிகோளமே மேலுள்ள நிறுவலைச்
சாத்தியமாக்குகிறது. அப்படியானால், சில பத்தாண்டுகளுக்கு முன்
மாற்றீட்டால் ``விரும்பத்தகாத'' விளைவுகள் எதுவும் இல்லை என்று
பூலோஸ் கூறியதையும் அவர் மறுபரிசீலனை செய்ய வேண்டியிருக்கும்
(\olref[replacement][extrinsic]{sec} ஐக் காண்க).

% SEGMENT OLP-0591-S03
ஆனால் $\kappa$ இருப்பது அவ்வளவு மோசமானதா? இங்கு ரசலின்
\emph{டிரிஸ்ட்ரம் ஷாண்டி முரண்புதிரை} நினைத்துப் பார்க்கலாம்.
டிரிஸ்ட்ரம் ஷாண்டி தன் வாழ்க்கையை நாட்குறிப்பில் பதிவு செய்கிறார்;
ஆனால் ஒரு நாளைப் பதிவுசெய்ய அவருக்கு ஓராண்டு ஆகிறது.
ஆண்டுகள் செல்லச் செல்ல அவர் மேலும் மேலும் பின்தங்குகிறார்:
முதலாண்டுக்குப் பின் ஒரு நாளையே பதிவுசெய்துள்ளார்; 364 நாட்கள்
பதிவாகவில்லை. இரண்டாம் ஆண்டுக்குப் பின் இரு நாட்கள் மட்டுமே
பதிவாகியுள்ளன; 728 நாட்கள் பதிவாகவில்லை. மூன்றாம் ஆண்டுக்குப்
பின் மூன்று நாட்கள் மட்டுமே பதிவாகியுள்ளன; 1092 நாட்கள்
பதிவாகவில்லை \dots\footnote{நெட்டாண்டுகளைப் பொருட்படுத்தவில்லை.}
இருப்பினும், டிரிஸ்ட்ரம் \emph{இறவாதவர்} எனில், அவர் ஒவ்வொரு
நாளையும் இறுதியில் பதிவுசெய்வார்; வாழ்வின் $n$ ஆம் ஆண்டில்
$n$ ஆம் நாளைப் பதிவுசெய்வார். ஆகவே ``காலத்தின் முடிவில்''
அவருடைய நாட்குறிப்பு முழுமையாகியிருக்கும்.

இது, $\alpha$ ஆனது $\aleph_\alpha$ ஐவிடச் சிறியதாக,
மேலும் மேலும் மிகவும் சிறியதாக, $\kappa$ வரை இருந்து,
அந்தப் புள்ளியில் $\kappa = \aleph_\kappa$ எனச் சமமாவதை
நினைப்பதிலிருந்து ஏன் அவ்வளவு வேறுபட வேண்டும்?

% SEGMENT OLP-0591-S04
அதை ஒருபுறம் வைத்து $\ZFC$ ஐ ஏற்றுக்கொள்வோம். நிலைப்புள்ளிக்
கட்டமைப்புகளைப் பற்றி இன்னும் சில முடிவுகளுடன் முடிப்போம்.
அடுத்த மூன்று முடிவுகள், படிநிலைமுறை அதன் உயரத்துக்கு
இணையான அகலத்தை அடையும் ஒரு முக்கியப் புள்ளி இருப்பதைக்
உள்ளுணர்வாகக் காட்டுகின்றன:

\begin{prop}\ollabel{bethfixed}
ஒரு $\beth$-நிலைப்புள்ளி உள்ளது; அதாவது $\kappa=
\beth_\kappa$ ஆகும் ஒரு $\kappa$ உள்ளது.
\end{prop}

\begin{proof}
\olref{alephfixed} இல் ``$\aleph$'' க்குப் பதில்
``$\beth$'' ஐப் பயன்படுத்திய அதே நிறுவல் பொருந்தும்.
\end{proof}

\begin{prop}\ollabel{stagesize}
$\card{V_{\omega+\alpha}} = \beth_{\alpha}$. மேலும்
$\omega \ordtimes \omega \leq \alpha$ எனில்
$\card{V_\alpha} = \beth_\alpha$.
\end{prop}

\begin{proof}
முதல் கூற்று எளிய முடிவிலிக்கடந்த தொகுத்தறிதலால் கிடைக்கிறது.
இரண்டாவது கூற்று, $\omega \ordtimes \omega \leq \alpha$
எனும்போது $\omega + \alpha = \alpha$ என்பதிலிருந்து கிடைக்கும்.
இதை நிரூபிக்க \olref[ord-arithmetic][]{chap} இன் வரிசையெண்
கணித முடிவுகளைப் பயன்படுத்துவோம். முதலில்
$\omega \ordtimes \omega = \omega \ordtimes (1 \ordplus \omega) =
(\omega  \ordtimes 1) \ordplus (\omega\ordtimes\omega) = \omega
\ordplus (\omega \ordtimes \omega)$ என்பதை கவனிக்கவும். இப்போது
$\omega \ordtimes \omega \leq \alpha$ எனில், ஏதேனும் $\beta$
க்கு $\alpha = (\omega\ordtimes\omega) \ordplus \beta$.
ஆகவே $\omega \ordplus \alpha = \omega \ordplus
((\omega \ordtimes \omega) \ordplus \beta) = (\omega \ordplus (\omega
\ordtimes \omega)) \ordplus \beta = (\omega \ordtimes \omega) \ordplus
\beta = \alpha$.
\end{proof}

\begin{cor}
$\card{V_\kappa} = \kappa$ ஆகும் ஒரு $\kappa$ உள்ளது.
\end{cor}

\begin{proof}
\olref{bethfixed} இல் கிடைத்த ஒரு $\beth$-நிலைப்புள்ளியை
$\kappa$ என்க. தெளிவாக $\omega \ordtimes \omega < \kappa$.
எனவே \olref{stagesize} இன்படி
$\card{V_\kappa} = \beth_\kappa= \kappa$.
\end{proof}

% SEGMENT OLP-0591-S05
$V_\kappa$ இன் கீழுள்ள படிநிலைகளின் எண்ணிக்கையும்
$V_\kappa$ இன் !!{element}களின் எண்ணிக்கையும் ஒன்றே.
உள்ளுணர்வாக, $V_\kappa$ இன் அகலமும் உயரமும் சமம்.
இது டிரிஸ்ட்ரம் ஷாண்டியின் கதையை நினைவூட்டுகிறது: ஒவ்வொரு
படிநிலையிலிருந்து அடுத்ததற்குச் செல்லும்போதும் நாம்
\emph{அடுக்குக் கணத்தை} எடுக்கிறோம்; இதனால் படிநிலைமுறை
ஒவ்வொரு முறையும் \emph{மிகவும்} பெரிதாகிறது. ஆனால்
``முடிவில்'', அதாவது $\kappa$ ஆம் படிநிலையில், அதன் அகலம்
அதன் உயரத்தை எட்டுகிறது.

% SEGMENT OLP-0591-S06
\emph{படிநிலைமுறையின் அகலம் எத்தனை முறை அதன் உயரத்துக்குச்
சமமாகிறது?} என்று கேட்கலாம். பதில்: \emph{வரிசையெண்கள்
எத்தனை உள்ளனவோ அத்தனை முறை.} இதற்குச் சிறிது விளக்கம் தேவை.

$\tau$ என்ற பதத்தைப் பின்வருமாறு வரையறுக்கிறோம். எந்த $A$
க்கும்:
% SOURCE CORRECTION TA-STH-038: தொடக்கப் பெத் குறியீட்டை card(A) இன் தொடர்ச்சிக்குச் செலுத்தி கண்டிப்பான பெருக்கத்தை உறுதிசெய்க.
\begin{align*}
	\tau_0(A) & \defis \beth_{\card{A} \ordplus 1}\\
	\tau_{n+1}(A) & \defis \beth_{\tau_n(A)}\\
	\tau(A) & \defis \bigcup_{n < \omega}\tau_n(A)
\intertext{\olref{bethfixed} இல் போலவே, $\tau(A)$ ஒரு
$\beth$-நிலைப்புள்ளி; இது ஒவ்வொரு $A$ க்கும் பொருந்தும். மேலும்
தொடக்கக் குறியீடு தொடர்ச்சியெண் என்பதால்
$\card{A} < \tau(A)$.
இப்போது பின்வரும் மீள்வரையறையைக் கருதுக:}
	W_0 &\defis 0\\
	W_{\alpha + 1} & \defis \tau(W_\alpha)\\
	W_\alpha & \defis \bigcup_{\beta < \alpha} W_\beta
	\text{, $\alpha$ ஓர் எல்லை வரிசையெண் எனில்}
\end{align*}
இந்தக் கட்டமைப்பு எல்லா வரிசையெண்களுக்கும் வரையறுக்கப்பட்டுள்ளது.
உள்ளுணர்வாக, $W$ என்பது வரிசையெண்களிலிருந்து
$\beth$-நிலைப்புள்ளிகளுக்கு ஒரு ``!!a{injection}'' ஆகும்.
முன்போலவே, எந்த $\alpha$ க்கும் $V_{W_\alpha}$ இன் அகலம்
அதன் உயரத்துக்குச் சமம்.

\end{document}
